\section{Limiting regimes}
Table~\ref{tab:limits} compares three limiting regimes. When both heat capacities equal $C$, the mean is the arithmetic mean and the rate is $2G/C$. Both temperatures evolve, which accounts for the factor of two relative to a body exchanging heat with a fixed-temperature bath of the same conductance.

\begin{table}[ht]
\caption{Limits of the stipulated two-reservoir model. The bath limit holds at fixed finite time, fixed $C_1$, fixed $G$, and fixed initial temperatures.}\label{tab:limits}
\begin{ruledtabular}
\begin{tabular}{lll}
Regime & $T_*$ & $\lambda$ \\
\hline
$C_1=C_2=C$ & $[T_1(0)+T_2(0)]/2$ & $2G/C$ \\
$C_2\to\infty$ & $T_2(0)$ & $G/C_1$ \\
$G\to0^+$ & unchanged & $0^+$
\end{tabular}
\end{ruledtabular}
\end{table}

At fixed finite time the limit $C_2\to\infty$ leaves the second temperature unchanged and recovers a single exponential for the first. Taking $G\to0^+$ instead makes the relaxation time diverge. At exactly $G=0$ every initial pair is stationary, so the common-temperature conclusion applies to positive conductance and long-time relaxation, not to zero conductance itself. The long-time conclusion therefore requires keeping this endpoint distinct from the limit of positive conductance.
